\subsection{子宫后位：正常变异，不是病}

体检或超声报告上"子宫后位"几个字，常让人心里一紧：是不是病？会影响怀孕吗？是性生活痛的原因吗？本节一次说清。子宫在盆腔内的位置分为前位、中位（水平位）与后位，\textbf{约五分之一的女性子宫呈后位}——这是正常的解剖变异，就像有人是左撇子，多数人天生如此；少数情况是后天形成的（分娩后韧带支撑暂时减弱、子宫内膜异位症或盆腔炎形成的粘连牵拉等）。

\vspace{0.5em}
\noindent\textbf{对性生活的影响}：

\begin{itemize}
	\item \textbf{后位本身通常不引起性交痛}：绝大多数后位女性没有任何症状；若出现\textbf{深部性交痛}，应排查的是子宫内膜异位症、子宫腺肌症、盆腔粘连、盆腔炎等病因（见"深层性交痛"一节），而不是"后位"本身背锅
	\item \textbf{少数严重后屈者}在深插入或某些角度下可能感到坠胀不适——调整体位与插入深度即可缓解，无需治疗
\end{itemize}

\vspace{0.5em}
\noindent\textbf{对受孕的影响}：

\begin{itemize}
	\item \textbf{"子宫后位难怀孕"是流传很广的误解}：宫颈的位置与精子进入宫颈口的关系，并不因前位/后位而有实质差别；目前没有可靠证据表明单纯后位会降低受孕率，也没有证据表明"同房后垫高臀部"能纠正后位、提高受孕率
	\item \textbf{真正需要关注的是伴随病因}：由子宫内膜异位症或盆腔炎粘连导致的\textbf{获得性}后位固定，可能同时影响受孕——需要治疗的是原发病，而非"把子宫摆正"；临床上也并无"子宫复位"的常规治疗
	\item 试孕半年以上未果（35 岁以上半年），按不孕评估流程检查即可（见本章"不孕症的常见原因"），不必纠结体位
\end{itemize}

\vspace{0.5em}
\noindent\textbf{什么时候需要就医}：后位\textbf{合并}进行性加重的痛经、深部性交痛、月经异常或不孕时，就诊排查内异症与盆腔粘连；无任何症状的体检发现，无需任何处理。

\begin{tcolorbox}[colback=pink!5!white,colframe=red!50!black,title=贴心小叮咛]
"子宫后位"是方位描述，不是诊断。看到报告上的这四个字，把它当作"单眼皮"一样的个体差异就好——该干嘛干嘛。
\end{tcolorbox}
